\begin{abstract}
  Accurate microclimate environmental measurements using low-cost transducers, such as the AHT10 capacitive sensor, are frequently compromised by inherent hardware noise and thermal anomalies. 
  This study investigates the metrological characterization of the measurement chain and the optimization of digital signal filtering utilizing the Exponential Moving Average (EMA) algorithm. 
  To eliminate external thermal coupling (Joule heating) and cross-interference originating from the ESP32 microcontroller's integrated radio frequency (RF) module, data acquisition was strictly executed via hardware-timer deterministic sampling at a constant interval of $\Delta t = 2.051$ seconds with all wireless protocols fully disabled.
  Steady-state characterization within an isolated chamber yielded a pure baseline noise floor with a standard deviation of $\sigma = 0.0108^\circ\text{C}$ post-linear detrending. 
  This measured deviation exceeds the sensor's 20-bit analog-to-digital theoretical resolution ($1 \text{ LSB} \approx 0.00019^\circ\text{C}$) by a factor of 56, strongly indicating that the accuracy's lower bound is intrinsically dominated by analog front-end thermodynamic noise rather than digital quantization errors.

  To objectively calibrate the physical trade-off between noise attenuation and algorithmic phase lag, a calculus-based Cost Function was modeled using balanced baseline penalty weights ($W_1 = W_2 = 0.5$). 
  The numerical computation converged on an optimal relative smoothing factor of $\alpha = 0.1494$. Quantitative evaluation of the transient response to a thermal step input demonstrated that this parameter establishes an algorithmic time constant of $\tau \approx 12.67$ seconds, 
  providing a theoretical white noise standard deviation reduction of $71.6\%$. 
  Furthermore, the optimized filter effectively suppressed transient deviations, yielding a Root Mean Square Error (RMSE) of $0.0771^\circ\text{C}$ and capping the Max Lag Error at $0.1333^\circ\text{C}$. 
  These findings demonstrate that a mathematically optimized EMA filter can robustly balance signal integrity with the sensor's physical response inertia without relying on heuristic visual approximations.
\end{abstract}